\paragraph{Proposition E.2.20 (Proved: the explicit smooth plumbing presentation).}
The following finite smooth plumbing presentation uses only standard
pieces and specified attaching maps. Its conversion into ordinary
framed handles and the relative smooth identification are given in E.2.21.

Let $F=\mathbb R^4/\mathbb Z^4$, with basis $(g,u,w,d)$, and use the
two matrices $A_j$ and fixed vectors $v_j$ of A.1.1.  Set
\[
 (m_1,m_2)=(3,4),\quad(\kappa_1,\kappa_2)=(1,-1),\quad
 w_j=\kappa_jv_j,\qquad L=\ker\gamma,
\]
so that $\Lambda=\mathbb Zw_j\oplus L$ and $A_jw_j=w_j$.
On $L$ the two matrices are
\[
 B_1=\begin{pmatrix}0&1&0\\-1&-1&0\\1&0&1\end{pmatrix},
 \qquad
 B_2=\begin{pmatrix}0&-1&0\\1&0&0\\0&1&1\end{pmatrix}.
 \tag{E.U1}
\]
Define the standard smooth four-manifolds
\[
 \mathcal S_j=(T^3\times\mathbb R)/
       ((l,\sigma+1)\sim(B_j^{-1}l,\sigma)),
 \qquad \mathcal N_j=\mathcal S_j\times D^2_r.
 \tag{E.U2}
\]
In particular the cap's oriented normal disk is part of the
definition, not an unspecified disk bundle.

The cap boundary is identified with $M_{A_j}$ by the following
explicit map.  Write $x=\phi w_j+l$ and use the interval marking
$(x,1)\sim(A_j^{-1}x,0)$ on $M_{A_j}$.  Then
\[
 [x,\theta]\longmapsto
 \left([l,\sigma=m_j\kappa_j\phi+\theta],
                         \ r=e^{2\pi i\kappa_j\phi}\right).
 \tag{E.U3}
\]
All angles on the right are taken in their stated mapping-torus
quotients.  In the ordered angular variables $(\phi,\theta)$
and $(\sigma,\alpha)$, where $r=e^{2\pi i\alpha}$, the map is
\[
 \begin{pmatrix}\sigma\\\alpha\end{pmatrix}
 =
 \begin{pmatrix}m_j\kappa_j&1\\\kappa_j&0\end{pmatrix}
 \begin{pmatrix}\phi\\\theta\end{pmatrix},
 \qquad
 \begin{pmatrix}\phi\\\theta\end{pmatrix}
 =
 \begin{pmatrix}0&\kappa_j\\1&-m_j\end{pmatrix}
 \begin{pmatrix}\sigma\\\alpha\end{pmatrix}.
 \tag{E.U4}
\]
These are integral unimodular matrices.

To define a compact six-manifold $\mathcal E_{\rm pl}$, take the
flat $F$-bundle over an oriented disc with three open interior
discs removed.  Its three marked inner loop monodromies are
$A_1,A_2,M_0$, with $A_1A_2M_0=I$.  Equivalently cut the planar
surface along three disjoint paths to its outer boundary, take
the product with $F$, and reglue the two banks by these three
integral linear diffeomorphisms.  This specifies smooth band
attachments, including their product normal coordinates.
Fill the first two boundary components by (E.U2)--(E.U4), and
the third by the finite canonical-bundle plumbing of E.2.19,
using its complete marked $M_0$ boundary map.  The remaining
boundary is marked as $F\times S^1_\mu$, with the outer meridian
orientation of (E.S1).  The core uses its zero section, so no
affine translation remains unspecified on that boundary.

For the actual original regular-fibre complement there is an
orientation-preserving diffeomorphism
\[
                  E\ \cong\ \mathcal E_{\rm pl}
                 \quad\hbox{relative to }F\times S^1_\mu .
 \tag{E.U5}
\]
Consequently the shifted fivefold has the explicit smooth
presentation
\[
 P\ \cong\
 (\mathcal E_{\rm pl}\times T^4)
 \ \cup_{\ (x_1,x_2,\theta)\mapsto(x_1,x_2,-\theta)\ }\
 (T^4\times\mathcal E_{\rm pl}).
 \tag{E.U6}
\]
Every piece in this presentation is a standard torus product,
a linear three-torus mapping torus times an oriented disk, or
the canonical disk bundle of the explicitly specified toric
blow-up with the normal-disk plumbings (E.T2).  All attachment
maps have been specified on their normal coordinates.

\begin{proof}
We verify first that (E.U2)--(E.U4) are the actual finite caps,
including the normal framing.  The original cap is
\[
 (F\times D^2_s)/
 \bigl((x,s)\sim(A_jx+v_j/m_j,e^{-2\pi i/m_j}s)\bigr).
 \tag{E.U7}
\]
The good-reduction real flat marking is smooth on the entire
root disc, so these are smooth coordinates even at $s=0$.
Because $\gamma A_j=\gamma$ and $\gamma(v_j)=\kappa_j=\pm1$,
\[
                  r=s\,e^{2\pi i\kappa_j\gamma(x)}
 \tag{E.U8}
\]
is invariant under the displayed action.  It is single-valued
on $F$, has the same absolute value as $s$, and is complex
linear on every normal disk.  The induced map
$[x,s]\mapsto([x],r)$ identifies (E.U7) smoothly with
$S_j\times D^2$.  An inverse is obtained by choosing any
representative $x$ and setting
$s=r e^{-2\pi i\kappa_j\gamma(x)}$; invariance makes this
independent of the representative in the quotient.

In the integral basis $(w_j,u,w,d)$ the affine generator is
\[
                  (\phi,l)\longmapsto
                     (\phi+\kappa_j/m_j,B_jl).
 \tag{E.U9}
\]
Thus $\sigma=m_j\kappa_j\phi$ gives exactly the mapping torus
in (E.U2).  This also supplies a direct inverse to the product
description: choose a representative $(l,\sigma)$, put
$x=(\kappa_j\sigma/m_j)w_j+l$ and
$s=r e^{-2\pi i\sigma/m_j}$.  Replacing that representative
by the defining mapping-torus relation changes $(x,s)$ by
the finite affine generator and hence not its class.

The actual boundary parametrization from (B.5.3) is
\[
 [x,\theta]\longmapsto
       [x+\theta v_j/m_j,e^{-2\pi i\theta/m_j}].
 \tag{E.U10}
\]
Substitute (E.U10) into (E.U8)--(E.U9).  The coefficient of
$\theta$ in the normal angle is
$-1/m_j+\kappa_j^2/m_j=0$, while the base angle is
$m_j\kappa_j\phi+\theta$.  This proves (E.U3) and (E.U4).
At $\theta=1$ that base angle advances by one, precisely
replacing $l$ by $B_j^{-1}l$; this checks the mapping-torus
seam.  A lattice change of $\phi$ changes the base angle by
$m_j\kappa_j$, and $B_j^{m_j}=I$, so it introduces no further
ambiguity.

The function of the cap that gives its original base parameter
is, in these smooth coordinates,
\[
                         q=r^{m_j}e^{-2\pi i\sigma}.
 \tag{E.U11}
\]
On the marked boundary it has angle
$m_j\alpha-\sigma=-\theta$, the original clockwise base
convention.  A positive normal meridian at fixed
$(l,\sigma)$ has $(\Delta\phi,\Delta\theta)=(\kappa_j,-m_j)$.
Since $\kappa_jw_j=v_j$, it is exactly $t_{v_j}\rho_j^{-m_j}$,
recovering the actual filling relation
$\rho_j^{m_j}=t_{v_j}$.  The normal circle is therefore not
silently exchanged for a once-around base loop.  The exact
integer checks of all these seam, framing and meridian
identities are in \path{code/maps_finite_cap_framings.py}.

Over the punctured planar core, the actual family has the
smooth real flat marking with linear monodromies $A_j$ and
its zero section.  The logarithmic finite gluing has exactly
the drift in (E.U10); after subtracting it, the remaining
translation has a single-valued real lift and is removed by
the collar isotopy already proved in B.5.1.  No nonzero
lattice winding is discarded.  At the cusp use the actual
relative smooth diffeomorphism of E.2.19.  Each comparison
has precisely the prescribed boundary value and can be
made product on a smaller collar.  They therefore glue to
the relative diffeomorphism (E.U5).  In particular the
normal framing of the unfilled ordinary fibre is unchanged.

Finally use the two hemispheres in the actual shifted
fibre-product decomposition of E.2.6.  On each hemisphere
the other factor is a regular marked torus bundle over a
disc, hence a product preserving its zero section and flat
marking.  Their common boundary has unchanged torus
coordinates and opposite normal angles.  Taking the product
of (E.U5) with those regular factors and gluing gives
(E.U6), with its induced complex orientation.  This is a
comparison of all smooth pieces and all boundary maps,
not an identification inferred from group or homology.
\end{proof}

E.2.21 explicitly converts the canonical-bundle self-plumbing and
all the exterior attachments into ordinary framed handles.
In the absolute planar-core description used here, the band part
contains three products
$T^4\times D^1\times D^1$ attached by the specified linear maps;
its product handles include index five.  Those high handles
cannot all disappear: if a compact six-manifold with this boundary
had only handles of index at most four, its dual relative
decomposition would make $\pi_1(\partial E)\to\pi_1(E)$ surjective,
contradicting (E.S1).  The displayed plumbing presentation is
therefore not being confused with a four-dimensional thickening
or a chain cancellation.
